% Part: turing-machines
% Chapter: machines-computations
% Section: introduction

\documentclass[../../../include/open-logic-section]{subfiles}

\begin{document}

\olfileid{tur}{mac}{int}
\olsection{پېژندنه}

دا څه معنا لري چې، مثلاً، له~$\Nat$ څخه~$\Nat$ ته يوه تابع
\emph{محاسبه کېدونکې} وي؟ د لومړنيو او تر ټولو نامتو ځوابونو څخه يو
دا دے چې تابع هغه وخت محاسبه کېدونکې ده چې د ټيورينګ ماشين په وسيله
محاسبه کېدے شي۔ دا مفهوم الن ټيورينګ په~1936 کښې وړاندې کړ۔ د ټيورينګ
ماشينونه د \emph{محاسبې د يو مدل} بېلګه ده---دوي د ``محاسبوي کړنلارې''
نظر د تعريفولو يوه رياضياتي کره لار ده۔ د دې کره معنا تر بحث لاندې
ده، خو په پراخه کچه منل شوې چې ټيورينګ ماشينونه د محاسبوي کړنلارو د
ټاکلو يوه لار ده۔ که څه هم د ``ټيورينګ ماشين'' اصطلاح د خوځنده پرزو
لرونکي فزيکي ماشين انځور ذهن ته راولي، په کره معنا ټيورينګ ماشين يو
سوچه رياضياتي جوړښت دے، او په همدې توګه د محاسبوي کړنلارې نظر مثالي
کوي۔ مثلاً، موږ د ټيورينګ ماشين د وخت يا حافظې پر اړتياوو هېڅ بنديز
نۀ ږدو: ټيورينګ ماشينونه يو څيز ان هغه وخت محاسبه کولے شي چې محاسبه
له هغو ټولو ذرو څخه زياته زېرمه يا زيات ګامونه وغواړي چې په کاينات
کښې شته۔

\begin{explain}
ښايي غوره وي چې ټيورينګ ماشين د يو ځانګړي ډول خيالي ميکانيزم د
پروګرام په توګه وګڼو۔ دا ميکانيزم له يوې \emph{پټې} او يو
\emph{لوست-ليک سر} څخه جوړ دے۔ زموږ د ټيورينګ ماشينونو په بڼه کښې
پټه په يوه لوري (ښي لوري ته) نامتناهي ده، او په \emph{مربعونو} وېشل
شوې، چې هر يو ئې د يوې متناهي \emph{الفبا} يوه نښه لرلے شي۔ دا ډول
الفباوې هر شمېر بېلابېلې نښې لرلے شي، خو موږ به تر ډېره په درېيو بسنه
کوو: $\TMendtape$، $\TMblank$ او~$\TMstroke$۔ کله چې ميکانيزم پيل
شي، پټه تشه وي (يعنې هر مربع د~$\TMblank$ نښه لري)، پرته له تر ټولو
کيڼ مربع څخه چې~$\TMendtape$ لري، او يو متناهي شمېر مربعونو څخه چې
\emph{ننوت} لري۔ په هر وخت کښې ميکانيزم له متناهي شمېر
\emph{حالتونو} څخه په يوه کښې وي۔ په پيل کښې سر د ننوت تر ټولو کيڼ
مربع، يعنې د پټې د پای نښې ښي لوري ته مربع، لولي او په يو ټاکلي
\emph{لومړني حالت} کښې وي۔ د ميکانيزم د چلېدو په هر ګام کښې د اوس
لوستل کېدونکي مربع منځپانګه، د ميکانيزم حالت او د ټيورينګ ماشين
پروګرام په ګډه ټاکي چې ورپسې څه کېږي۔ د ټيورينګ ماشين پروګرام د يوې
جزوي تابع په توګه ورکول کېږي چې حالت~$q$ او نښه~$\sigma$ د ننوت په
توګه اخلي، او درې‌يز~$\tuple{q', \sigma', D}$ راګرځوي۔ هر کله چې
ميکانيزم په حالت~$q$ کښې وي او نښه~$\sigma$ لولي، د اوسني مربع نښه
په~$\sigma'$ بدلوي، سر د دې له مخې کيڼ، ښي لوري ته ځي يا پر ځاے پاتې
کېږي چې~$D$ په ترتيب سره~$\TMleft$، $\TMright$ يا~$\TMstay$ وي، او
ميکانيزم حالت~$q'$ ته ځي۔
\textbf{د ژباړې د نسخې سمون (OLCMP-039):} سرچينې د پټې تر ټولو کيڼ
مربع د پای نښې د ځاے په توګه ټاکلے و، خو بيا ئې سر هم په پيل کښې پر
هماغه مربع ايښے و۔ د همدې څپرکي صوري لومړنی ترتيب سر د پای نښې ښي
لوري ته، د ننوت پر لومړي مربع ږدي؛ دلته هماغه ټاکلے دود په څرګنده
ساتل شوے دے۔

د بېلګې په توګه، په~\olref{fig:tm} کښې حالت وګورئ۔
\begin{figure}
  \olasset{\olpath/assets/diagrams/turing-machine.tikz}
  \caption{يو ټيورينګ ماشين چې خپل پروګرام عملي کوي۔}
  \ollabel{fig:tm}
\end{figure}
د ټيورينګ ماشين د پټې ليدونکې برخه په تر ټولو کيڼ مربع کښې د پټې د
پای نښه~$\TMendtape$ لري، ورپسې درې~$1$ ګانې، يو~$0$، او بيا څلور
نورې~$1$ ګانې دي۔ سر له کيڼ لوري درېيم مربع لولي، چې~$1$ لري، او په
حالت~$q_1$ کښې دے---وايو ``ماشين په حالت~$q_1$ کښې~$1$ لولي۔'' که
د ټيورينګ ماشين پروګرام د ننوت~$\tuple{q_1, 1}$ دپاره درې‌يز
$\tuple{q_2, 0, \TMstay}$ را وګرځوي، ماشين به اوس په درېيم مربع کښې
~$1$ په~$0$ بدل کړي، لوست-ليک سر به پر خپل ځاے پرېږدي، او حالت~$q_2$
ته به واوړي۔ که بيا پروګرام د ننوت~$\tuple{q_2, 0}$ دپاره
$\tuple{q_3, 0, \TMright}$ را وګرځوي، ماشين به~$0$ په بل~$0$ وليکي
(يعنې په عمل کښې به د لوست-ليک سر لاندې د پټې منځپانګه هماغسې
پرېږدي)، يو مربع ښي لوري ته به لاړ شي، او حالت~$q_3$ ته به ننوځي۔ او
همداسې نور۔

وايو ماشين هغه وخت \emph{درېږي} چې له داسې حالت~$q_n$ او نښې
~$\sigma$ سره مخ شي چې د~$\tuple{q_n, \sigma}$ دپاره هېڅ لارښوونه
نۀ وي، يعنې د ننوت~$\tuple{q_n,\sigma}$ دپاره د حالت بدلون تابع
تعريف شوې نۀ وي۔ په بله وينا، ماشين د عملي کولو دپاره هېڅ لارښوونه
نۀ لري، او په هماغه ځاے کار بندوي۔ کله ناکله درېدل د ځانګړي درېدني
حالت~$h$ په وسيله څرګندېږي۔ دا به وروسته په زياتو جزئياتو وښودل شي۔
\end{explain}

\begin{digress}
د ټيورينګ د مقالې، ``د محاسبه کېدونکو عددونو په اړه،'' ښکلا دا ده چې
هغه نۀ يوازې يو صوري تعريف وړاندې کوي، بلکې دا استدلال هم کوي چې
تعريف د محاسبه کېدنې شهودي مفهوم رانغاړي۔ له تعريفه بايد څرګنده وي
چې هره تابع چې د ټيورينګ ماشين په وسيله محاسبه کېږي، په شهودي معنا
محاسبه کېدونکې ده۔ ټيورينګ درې ډوله استدلالونه وړاندې کوي چې معکوسه
خبره هم رښتيا ده، يعنې هره تابع چې موږ ئې په طبيعي ډول محاسبه کېدونکې
ګڼو، د داسې ماشين په وسيله محاسبه کېدونکې ده۔ د ټيورينګ په خپلو
ټکو، دا دي:
\begin{enumerate}
\item شهود ته نېغه مراجعه۔
\item د دوو تعريفونو د معادلتوب ثبوت (په هغه صورت کښې چې نوی تعريف
  زيات شهودي جاذبيت ولري)۔
\item د محاسبه کېدونکو عددونو د لويو ټولګيو بېلګې ورکول۔
\end{enumerate}
زموږ موخه دا ده چې د محاسبه کېدنې مفهوم ``په اصل کښې'' تعريف کړو،
يعنې د وخت او ځاے عملي محدوديتونه په پام کښې وانۀ خلو۔ البته، کله چې
د محاسبه کېدنې تر ټولو پراخ تعريف ولرو، بيا د محدودو سرچينو لرونکې
محاسبه څېړلے شو؛ دا د ``محاسبوي پېچلتيا'' په نوم مضمون زړه جوړوي۔
\end{digress}

\begin{history}
الن ټيورينګ په~1936 کښې ټيورينګ ماشينونه اختراع کړل۔ که څه هم په هغه
وخت کښې د هغه پاملرنه د لومړۍ درجې منطق پرېکړه کېدنې ته وه، دا مقاله
د کمپيوټر د جوړښت د بنسټونو يوه پرېکنده مقاله بلل شوې ده۔ په مقاله
کښې ټيورينګ پر محاسبه کېدونکو حقيقي عددونو، يعنې هغو حقيقي عددونو چې
اعشاري غځونې ئې محاسبه کېدونکې دي، تمرکز کوي؛ خو يادونه کوي چې د هغه
مفهومونه پر طبيعي عددونو محاسبه کېدونکو تابعو او نورو ته اړول ستونزمن
نۀ دي۔ پام وکړئ چې دا په~1941 کښې د لومړي کارکوونکي عمومي-موخې
کمپيوټر له جوړېدو پوره پنځه کاله مخکښې وو (الماني کونراد زوزه دا د
خپلو والدينو د کور په ناسته کوټه کښې جوړ کړ)، په بلېچلي پارک کښې د
ټيورينګ او د هغه د ملګرو د کوډ ماتوونکي کولوسس له جوړېدو اووه کاله
مخکښې (1943)، د امريکايي اېنياک څخه نهه کاله مخکښې (1945)، په
مانچسټر کښې د لومړي برتانوي عمومي-موخې کمپيوټر---د مانچسټر د کوچنۍ
کچې ازمېښتي ماشين---له جوړېدو دولس کاله مخکښې (1948)، او د
امريکايانو له خوا د بيناک له لومړۍ ازموينې ديارلس کاله مخکښې (1949)۔
د مانچسټر کوچنۍ کچې ازمېښتی ماشين د لومړي زېرمه-پروګرام کمپيوټر وياړ
لري---پخواني ماشينونه د هرې نوې دندې دپاره په لاس بيا سيمبندي کېدل۔
\end{history}

\end{document}
